\chapter{Đối Đáp Trong Giờ Sinh Hoạt}
\chapterintro{Giờ sinh hoạt cần vừa mềm vừa rõ}{Đây là nơi một câu có duyên thường mở được lòng lớp hơn một bài nói dài.}

\section{Tuần này lớp mình đỡ nhất ở điểm nào?}
\begin{manthoai}{Màn thoại}
\textbf{Thầy/Cô:} Tuần này lớp mình đỡ nhất ở điểm nào?\\
\textbf{Học sinh:} Đỡ trễ hơn ạ.\\
\textbf{Thầy/Cô:} Tốt. Mình bắt đầu từ tiến bộ trước đã.
\end{manthoai}
\begin{dungkhinào}
Rất hợp để mở giờ sinh hoạt bằng năng lượng tích cực.
\end{dungkhinào}
\begin{nhipthem}
Cho lớp nói điểm tốt trước, rồi mới sang điều cần sửa.
\end{nhipthem}
\ngancach

\section{Lớp mình đang cần bớt drama hay bớt lười?}
\begin{manthoai}{Màn thoại}
\textbf{Thầy/Cô:} Lớp mình đang cần bớt drama hay bớt lười?\\
\textbf{Học sinh:} Cả hai ạ.\\
\textbf{Thầy/Cô:} Rồi, chọn một cái sửa trước thôi.
\end{manthoai}
\begin{dungkhinào}
Hợp để nhìn thẳng vào vấn đề mà vẫn giữ được tiếng cười nhẹ.
\end{dungkhinào}
\begin{nhipthem}
Sau câu vui, chốt ngay một việc cụ thể cho tuần tới.
\end{nhipthem}
\ngancach

\section{Ai thấy mình cần sửa ít nhất một điều thì giơ tay thật}
\begin{manthoai}{Màn thoại}
\textbf{Thầy/Cô:} Ai thấy mình cần sửa ít nhất một điều thì giơ tay thật.\\
\textbf{Học sinh:} Cả lớp giơ tay.\\
\textbf{Thầy/Cô:} Tốt. Vậy mình nói chuyện như người cùng sửa, không phải người đi bắt lỗi.
\end{manthoai}
\begin{dungkhinào}
Rất hợp để hạ không khí phòng thủ của lớp.
\end{dungkhinào}
\begin{nhipthem}
Sau đó mời 2 đến 3 bạn nói điều mình muốn sửa trong tuần mới.
\end{nhipthem}
\ngancach

\section{Tuần này lớp muốn được nhớ vì điều gì?}
\begin{manthoai}{Màn thoại}
\textbf{Thầy/Cô:} Tuần này lớp muốn được nhớ vì điều gì?\\
\textbf{Học sinh:} Vì đỡ ồn hơn ạ.\\
\textbf{Thầy/Cô:} Rồi, vậy mình làm sao để lời giới thiệu đó là thật.
\end{manthoai}
\begin{dungkhinào}
Hợp khi muốn giờ sinh hoạt có thêm chất xây dựng hình ảnh tập thể, không chỉ xử lý lỗi.
\end{dungkhinào}
\begin{nhipthem}
Câu này mở ra mục tiêu chung rất tốt cho tuần kế tiếp.
\end{nhipthem}
\ngancach

\section{Lớp hứa với nhau một điều ít thôi, nhưng làm thật}
\begin{manthoai}{Màn thoại}
\textbf{Thầy/Cô:} Lớp hứa với nhau một điều ít thôi, nhưng làm thật, là điều gì?\\
\textbf{Học sinh:} Không vào lớp muộn ạ.\\
\textbf{Thầy/Cô:} Ít mà thật thì mạnh hơn nhiều lời hứa đẹp.
\end{manthoai}
\begin{dungkhinào}
Dùng để chốt cuối buổi sinh hoạt khi cần một cam kết ngắn và rõ.
\end{dungkhinào}
\begin{nhipthem}
Chỉ nên lấy đúng một việc tập thể có thể tự nhắc nhau trong tuần sau.
\end{nhipthem}
